% 胖打散的间隔带：X=N，f_C(i)=1/i (i∈C) 或 0 (i∉C)。
% 每个点 i 的两个可选值为 0 与 1/i，阈值取中点 r_i=1/(2i)，尺度 gamma=0.12。
% 间隔带 [r_i-gamma, r_i+gamma] 必须夹在两个可选值之间，即 1/(2i)>=gamma，i<=1/(2gamma)=4.17。
% 因而 i=1,...,4 可在该尺度下同时被胖打散，i=5,6 不能；fat_{0.12}=4。
% 纵轴比例：值 1 对应 40 mm；横向间距 15 mm。
\begingroup
\providecolor{FigureGreen}{HTML}{829C8B}
\providecolor{FigureGreenFill}{HTML}{E8EEE9}
\begin{tikzpicture}[figure picture,x=1mm,y=1mm]
\def\g{0.12}
\draw[FigureStroke,line width=.6pt,-{Stealth[length=1.6mm]}] (0,0)--(100,0) node[right]{$i$};
\draw[FigureStroke,line width=.6pt,-{Stealth[length=1.6mm]}] (0,0)--(0,45) node[above]{$f(i)$};
\foreach \y/\t in {0/0,0.25/0.25,0.5/0.5,1/1}{\draw[FigureStroke] (-1,{40*\y})--(0,{40*\y}); \node[anchor=east,figure note] at (-1.5,{40*\y}) {$\t$};}
\foreach \i in {1,...,6}{
 \pgfmathsetmacro{\top}{1/\i}
 \pgfmathsetmacro{\r}{1/(2*\i)}
 \pgfmathsetmacro{\lo}{\r-\g}
 \pgfmathsetmacro{\hi}{\r+\g}
 \pgfmathsetmacro{\ok}{\r>=\g ? 1 : 0}
 \pgfmathsetmacro{\xx}{15*\i}
 \ifdim\ok pt>0.5pt
  \fill[FigureGreenFill] ({\xx-4},{40*\lo}) rectangle ({\xx+4},{40*\hi});
  \draw[FigureGreen,line width=.6pt] ({\xx-4},{40*\lo}) rectangle ({\xx+4},{40*\hi});
 \else
  \fill[FigureRedFill] ({\xx-4},{40*max(\lo,-0.06)}) rectangle ({\xx+4},{40*\hi});
  \draw[FigureRed,line width=.6pt] ({\xx-4},{40*max(\lo,-0.06)}) rectangle ({\xx+4},{40*\hi});
 \fi
 \draw[FigureStroke,line width=.8pt,dash pattern=on 1.5pt off 1pt] ({\xx-5},{40*\r})--({\xx+5},{40*\r});
 \draw[FigureMuted,line width=.5pt] (\xx,0)--(\xx,{40*\top});
 \fill[FigureBlue] (\xx,{40*\top}) circle (1.1);
 \fill[FigureBlue] (\xx,0) circle (1.1);
 \node[figure note] at (\xx,-4) {$\i$};
}
\node[figure note,anchor=west] at (22,42) {圆点：可选值 $0$ 与 $1/i$；虚线：阈值 $r_i=1/(2i)$};
\node[figure note,anchor=west] at (22,37) {色带：$[r_i-\gamma,\,r_i+\gamma]$，$\gamma=0.12$};
\node[figure note,text=FigureGreen] at (37.5,-9) {可胖打散};
\node[figure note,text=FigureRed] at (82.5,-9) {带宽超过两值间距};
\end{tikzpicture}
\endgroup
